\scanpage{119}
\mspage{3}
% archive sheet number: 118 (upper left)
donnée \deleted{\unsure{est}} \inserted{étant} soumise aux axiomes qui suivent $\mathrm{Cart}(1)$ à $\mathrm{Cart}(\ldots)$.

Comme je ne suis pas sûr de connaître d'avance \inserted{toutes} les propriétés essentielles, le nombre de ces axiomes reste pour le moment ouvert.

\uline{Cart(1)} (\uline{Symétrie}) Si un carré
\diagram{\begin{tikzcd}[ampersand replacement=\&]
X_1 \arrow[d] \& X_3 \arrow[l] \arrow[d] \\
X_0 \& X_2 \arrow[l]
\end{tikzcd}}
est \inserted{$h$-}cartésien, de $=$ le carré symétrique (obtenu en échangeant le rôle de $X_1$ et de $X_2$).

\uline{Cart(2)} (\uline{Carrés triviaux}) Soit
\diagram{\begin{tikzcd}[ampersand replacement=\&]
X_1 \arrow[d,"p_1"'] \& X_3 \arrow[l,"q_2"'] \arrow[d,"q_1"] \\
X_0 \& X_2 \arrow[l,"p_2"]
\end{tikzcd}}
un carré dans $\unsure{\mathscr H|\mathcal H|H}$, avec $p_1$ un iso. Alors le carré est $h$-cartésien \unsure{sss|ssi} $q_1$ est un iso.

\uline{Cart(3)} (\uline{Carrés composés}) Soit un diagr. comm.
\diagram{\begin{tikzcd}[ampersand replacement=\&]
X'_0 \arrow[d] \arrow[dr,phantom,"{\underline{\overline{\mathrm I}}}" description] \& X'_1 \arrow[l] \arrow[d] \arrow[dr,phantom,"{\underline{\overline{\mathrm{II}}}}" description] \& X'_2 \arrow[l] \arrow[d] \\
X_0 \& X_1 \arrow[l] \& X_2 \arrow[l]
\end{tikzcd}}
, d'où les carrés $\underline{\overline{\mathrm I}}$, $\underline{\overline{\mathrm{II}}}$ et
\diagram{\begin{tikzcd}[ampersand replacement=\&]
\underline{\overline{\mathrm{III}}}: \& X'_0 \arrow[d] \& X'_2 \arrow[l] \arrow[d] \\
\& X_0 \& X_2 \arrow[l]
\end{tikzcd}}
(qu'on appelle \uline{carré composé} (``horizontal'') des carrés $\underline{\overline{\mathrm I}}$ et $\underline{\overline{\mathrm{II}}}$). Supposons $\underline{\overline{\mathrm I}}$ $h$-cartésien. Alors $\underline{\overline{\mathrm{II}}}$ est $h$-cartésien sss $\underline{\overline{\mathrm{III}}}$ l'est.

\deleted{\uline{\unsure{Cart(4)}}} \inserted{Corollaire} \deleted{(\unsure{Théorème} \unsure{des cubes})} Soit un diagramme commutatif cubique, \inserted{\unsure{qu'on}} peut considérer comme hom. d'un carré \inserted{(supérieur du cube)} dans un autre (face inférieure) :
\diagram{\begin{tikzcd}[ampersand replacement=\&,row sep=small,column sep=small]
\& {} \arrow[dl] \arrow[dd,dashed] \& \& {} \arrow[ll] \arrow[dl] \arrow[dd] \\
{} \arrow[dd] \& \& {} \arrow[ll] \arrow[dd] \& \\
\& {} \arrow[dl] \& \& {} \arrow[ll,dashed] \arrow[dl] \\
{} \& \& {} \arrow[ll] \&
\end{tikzcd}}
Toutes les faces du cube\textsuperscript{*}, \uline{sf} la face supérieure sont
% margin note written slanted, rising to the right
\margin{* En fait, il suffit de la face inférieure, et de deux faces opposées latérales (p. ex. celle de \unsure{devant} et celle de derrière)}
